\chapter{Scoring compuesto}
\label{sec:paper-score}

\section{Motivación}

El DFS temporal produce paths estructuralmente válidos. Ninguno
está \emph{etiquetado} como malicioso. El scoring convierte el
conjunto de paths en una lista ordenada que un analista puede
triar. El diseño persigue tres propiedades:

\begin{enumerate}
    \item \textbf{Descomposable}. Un score único monolítico (``la
    caja negra dijo 0.87'') no informa al analista. Queremos que
    cada path muestre sus componentes y que el analista pueda
    razonar sobre ellos.
    \item \textbf{Operable}. Los pesos deben ser ajustables por
    entorno sin retraining de un modelo. Un analista con contexto
    organizacional sabe que \emph{novelty temporal} importa más
    en una telco que en una boutique financiera.
    \item \textbf{Degradable}. Si un componente no está disponible
    (típicamente \code{gnn_threat} sin modelo), el score total
    debe seguir siendo significativo.
\end{enumerate}

\section{Los cinco componentes}

Todos los componentes retornan $[0, 1]$. Sean $v$ un nodo, $e$ una
arista, $\pi = (v_0, e_0, \dots, v_n)$ un path.

\subsection{Entity rarity}

\[
\text{ER}(v) = 1 - \frac{\ln(\text{freq}(v))}{\ln(\max_u \text{freq}(u))}
\]

donde $\text{freq}(v)$ es el grado del nodo en el grafo actual.
El log normaliza la distribución de grados power-law. Rarity alta
para nodos que aparecen poco; rarity 0 para el hub dominante.
Aplicado por path como promedio sobre $v_0, \dots, v_n$.

\subsection{Edge rarity}

Análogo al anterior pero sobre pares $(\text{type}(v), \rho, \text{type}(u))$:

\[
\text{EdgeR}(e) = 1 - \frac{\ln(\text{freq}(\text{type}(e)))}{\ln(\max_{\rho'} \text{freq}(\rho'))}
\]

Captura qué tan raro es \emph{este tipo de relación} en el grafo.
Un \code{User-[Spawn]->Process} dentro de un grafo de red pura es
maxR; dentro de un grafo dominado por endpoint, es común.

\subsection{Neighborhood concentration}

\[
\text{NC}(\pi) = 1 - \frac{|\bigcap_{v \in \pi} N(v)|}{|\bigcup_{v \in \pi} N(v)|}
\]

Jaccard-complement entre los vecindarios de los nodos del path. Si
los vecindarios se solapan mucho (actividad rutinaria en la misma
sub-red), NC es bajo. Si los nodos del path se tocan pero
pertenecen a mundos distintos (lateral movement real), NC es alto.

\subsection{Temporal novelty}

\[
\text{TN}(\pi) = \min_{e \in \pi} \text{novelty}(e), \quad
\text{novelty}(e) = 1 - \frac{\text{n\_past}(e)}{\text{n\_past}(e) + 1}
\]

donde $\text{n\_past}(e)$ es la cantidad de veces que se observó
el mismo par ordenado $(v_u, v_v)$ en los 30 días previos. Una
arista nunca vista antes tiene \text{novelty} = 1; una arista
observada ayer 10 veces tiene $\sim 0.09$. El path toma el mínimo
porque una sola arista rutinaria basta para hacer el path poco
novedoso.

\subsection{GNN threat}

\[
\text{GT}(\pi) = \begin{cases} f_\theta(\pi) & \text{si hay modelo cargado}\\ 0 & \text{caso contrario} \end{cases}
\]

donde $f_\theta$ es un clasificador GNN entrenado sobre paths
etiquetados. El ``0 si no hay modelo'' no penaliza: el
compositor combina con peso $w_5 = 0$ por defecto (ver
cap.~\ref{sec:paper-gate}).

\section{Composición}

\[
S(\pi) = w_1\,\overline{\text{ER}} + w_2\,\overline{\text{EdgeR}} + w_3\,\text{NC} + w_4\,\text{TN} + w_5\,\text{GT}
\]

con $\sum w_i = 1$ y pesos por defecto
$(0.25, 0.25, 0.30, 0.20, 0.00)$. La ausencia de peso para GNN
obedece al principio de \emph{degradación segura}: el sistema
debe funcionar en el primer día sin modelo entrenado. Si el
cliente entrena uno, sube $w_5$ a expensas de los otros.

\section{Propiedades numéricas}

\begin{invariant}
$S(\pi) \in [0, 1]$ por construcción ya que cada componente
$\in [0,1]$ y $\sum w_i = 1$.
\end{invariant}

\begin{invariant}
Si dos paths $\pi, \pi'$ difieren sólo en un nodo $v \ne v'$ con
$\text{freq}(v) > \text{freq}(v')$, entonces
$\text{ER}(\pi) \le \text{ER}(\pi')$. \emph{Ceteris paribus},
favorece el path con entidades más raras.
\end{invariant}

El test \filepath{core/graph-engine/tests/anomaly_proptest.rs}
verifica estabilidad numérica: pequeñas perturbaciones en el
grafo no causan saltos en el score superiores a la perturbación.

\section{Explicación por path}

El endpoint \code{explain\_hunt\_result} devuelve los cinco
componentes normalizados y su producto con el peso. La UI los
presenta como barras apiladas. El diseño explicable es una
respuesta deliberada a la trampa de los scores opacos: un analista
que no entiende por qué un path subió no puede ni confiar ni
cuestionar el ranking.

La especificación formal de endpoints y fórmulas vive en el libro
\emph{Especificación}.
